\section{Two moving boundaries}
\label{sec:moving-cutoff}
%======================================================================

The decomposition changes with the cutoff.  The goal here is to represent that changing arithmetic classification as motion of two simple boundaries through one fixed point cloud.

Fix $R>0$.  The complete arithmetic prefix is
\begin{equation}
 \mathscr P_R
 :=\{w=x+iy\in\mathscr W:x<R^2\}.
 \label{eq:prefix-region}
\end{equation}
Since $V=q^2=|w|+y$, the smooth condition $q\le R$ is
\begin{equation}
 |w|+y\le R^2.
 \label{eq:smooth-radial}
\end{equation}
Define
\begin{align}
 \mathscr A_R
 &:={\{w\in\mathscr P_R:|w|+\Impart w\le R^2\}},
 \label{eq:A-region}\\
 \mathscr U_R
 &:={\{w\in\mathscr P_R:|w|+\Impart w>R^2\}}.
 \label{eq:U-region}
\end{align}
The prefix is the disjoint union
\begin{equation}
 \mathscr P_R=\mathscr A_R\,\dot\cup\,\mathscr U_R.
 \label{eq:region-partition}
\end{equation}

Solving \cref{eq:smooth-radial} on the boundary gives
\begin{equation}
 \boxed{
 \Gamma_R(x)=\frac{R^4-x^2}{2R^2}.
 }
 \label{eq:Gamma-R}
\end{equation}
Inside the prefix strip $x<R^2$, the smooth region lies below the graph and the transport region lies above it.

In null coordinates, the same regions are even simpler:
\begin{align}
 \mathscr P_R&:\quad UV<R^4,
 \label{eq:prefix-UV}\\
 \mathscr A_R&:\quad UV<R^4,\quad V\le R^2,
 \label{eq:smooth-UV}\\
 \mathscr U_R&:\quad UV<R^4,\quad V>R^2.
 \label{eq:transport-UV}
\end{align}
For a transport point, $UV<R^4<V^2$, so automatically $U<R^2$.  Thus
\begin{equation}
 \mathscr U_R:\quad U<R^2<V,
 \qquad
 UV<R^4.
 \label{eq:transport-null-clean}
\end{equation}
This is exactly $c<R<q$.

\subsection{Intersection with every cofactor curve}

\begin{proposition}[Intersection factorization]
\label{prop:intersection-factorization}
For $c,R>0$,
\begin{equation}
 \gamma_c(x)-\Gamma_R(x)
 =
 \frac{(R^2+c^2)(x^2-c^2R^2)}{2c^2R^2}.
 \label{eq:curve-gap}
\end{equation}
Consequently,
\begin{equation}
 \gamma_c(x)=\Gamma_R(x)
 \iff x=cR.
 \label{eq:curve-intersection}
\end{equation}
At the arithmetic point $x=cq$,
\begin{equation}
 \gamma_c(cq)-\Gamma_R(cq)
 =
 \frac{(R^2+c^2)(q^2-R^2)}{2R^2}.
 \label{eq:point-gap}
\end{equation}
\end{proposition}

The sign is therefore exact:
\begin{align}
 q<R&\iff w_{c,q}\text{ lies below }\Gamma_R,
 \label{eq:below}\\
 q>R&\iff w_{c,q}\text{ lies above }\Gamma_R.
 \label{eq:above}
\end{align}

\begin{proposition}[Orthogonal crossing]
\label{prop:orthogonal-crossing}
At the intersection $x=cR$,
\begin{equation}
 \gamma_c'(cR)=\frac{R}{c},
 \qquad
 \Gamma_R'(cR)=-\frac{c}{R},
 \label{eq:orthogonal-slopes}
\end{equation}
so
\begin{equation}
 \gamma_c'(cR)\Gamma_R'(cR)=-1.
 \label{eq:orthogonal-product}
\end{equation}
Hence each cofactor curve meets the moving cutoff orthogonally.
\end{proposition}

The transition is therefore transverse, never tangent.  A canonical point cannot remain trapped in a thin boundary layer as the cutoff passes its largest prime.

%======================================================================
